% ============================================================================
% Fichier  : partie1/P01_C01_investigation_numerique.tex
% Titre    : L'Investigation Numérique
% Auteur   : MINKA MI NGUIDJOÏ Thierry Emmanuel
% Version  : 1.0 -- novembre 2026
% Corrigés : corriges/Corr_C01.tex
% ============================================================================
% Sources : ISO/IEC 27037:2012 (principes, paraphrasés), ISO/IEC 27043:2015
% (classes de processus, paraphrasées), NIST SP 800-86 (quatre phases),
% NIST SP 800-61 r2, Locard (1920), Casey (2004), Heuer (1999, analyse des
% hypothèses concurrentes).
% Faits nouveaux (bible) : hypothèses H1 à H4 de la cellule du 15/03/2025.
% ============================================================================

\chapter{L'Investigation Numérique}
\label{chap:investigation}
\chaptermark{L'investigation numérique}

\epigraph{%
    « Une enquête ne cherche pas à confirmer une idée.\\
    Elle cherche ce qui pourrait la réfuter. »%
}{Principe du cours}

\niveauChapitre{Fondamental}{%
    Lecture du Chapitre~\ref{chap:prologue}
}

\begin{boxobjectifs}
\begin{itemize}
    \item Définir l'investigation numérique, la trace et la preuve
          numériques.
    \item Énoncer les principes qui rendent un examen acceptable~:
          traçabilité, répétabilité, reproductibilité, justification.
    \item Décrire les phases d'une enquête selon le NIST SP~800-86 et les
          classes de processus de l'ISO/IEC~27043.
    \item Distinguer enquête judiciaire, enquête interne et réponse à
          incident, et leurs exigences de preuve.
    \item Conduire un raisonnement par hypothèses concurrentes et coter
          chaque hypothèse.
\end{itemize}
\end{boxobjectifs}

\bigskip

% ============================================================================
\section*{Avant de commencer : quatre hypothèses sur la table}
% ============================================================================

Bafoussam, samedi 15 mars 2025, 8~heures. Mme~MBARGA, de SANKOFA EXPERTISE,
ouvre la première réunion de la mission en demandant à chacun d'écrire, sur
une feuille, comment il pense que l'attaquant est entré. Elle recueille
quatre réponses.

\begin{itemize}
    \item H1~: par la passerelle d'accès à distance, avec un compte volé ou
          une faille.
    \item H2~: par un poste d'agence infecté par un courriel piégé, puis par
          le réseau interne.
    \item H3~: par un employé du siège qui aurait agi de l'intérieur,
          depuis un poste du siège.
    \item H4~: par le prestataire de l'interface de paiement, dont la clé
          aurait fuité de son côté.
\end{itemize}

Le DSI tient pour H3, la rumeur du vendredi. Le directeur général tient pour
H4, qui met l'établissement hors de cause. Mme~MBARGA ne tient pour aucune~:
elle écrit les quatre au tableau et annonce que l'enquête consistera à
chercher, pour chacune, la pièce qui la réfuterait.

\begin{boxalert}
\textbf{Violation typique.} Choisir l'hypothèse la plus probable dès le
premier jour, puis ne chercher que les traces qui la confirment.

\smallskip
\textbf{Règle~:} on garde ouvertes toutes les hypothèses compatibles avec les
constatations, et l'on cherche en priorité la pièce capable d'en éliminer
une. Une hypothèse survit parce qu'on n'a pas réussi à la réfuter, pas parce
qu'on a trouvé des indices en sa faveur.
\end{boxalert}

\separateur

% ============================================================================
\section{Trace, Preuve et Investigation}
\label{sec:c01-definitions}
% ============================================================================

\begin{boxdef}
\textbf{\termecle{Trace numérique}}~: toute donnée produite par un système
d'information, conservée sur un support ou en transit, qui garde la marque
d'une action ou d'un état passé~: une ligne de journal, une date de
modification, un fichier supprimé mais récupérable, une connexion en mémoire.

\smallskip
\textbf{\termecle{Preuve numérique}}~: trace recueillie, préservée et
examinée de façon à pouvoir être présentée à un décideur (un juge, une
direction, un régulateur) et soutenir ou réfuter une proposition de fait.

\smallskip
\textbf{\termecle{Investigation numérique}}~: démarche méthodique qui
identifie, recueille, préserve, examine et analyse des traces numériques
pour répondre à des questions de fait, et en rend compte de façon
vérifiable. On parle aussi de \textit{digital forensics}.
\end{boxdef}

Toute trace n'est pas une preuve. Une trace devient preuve quand trois
conditions sont réunies~: on sait d'où elle vient, on peut montrer qu'elle
n'a pas changé depuis son recueil, et on peut expliquer ce qu'elle signifie.
Le Chapitre~\ref{chap:preuve} traite des deux premières~; la troisième
occupe toute la partie~II.

\subsection{Le principe d'échange}

Edmond Locard, fondateur du laboratoire de police technique de Lyon, a
formulé au début du \textsc{xx}\textsuperscript{e}~siècle l'idée que tout
contact laisse une trace~\cite{locard1920}. Le numérique la vérifie d'une
manière particulière. Un attaquant qui ouvre une session laisse une entrée
dans le journal de la passerelle~; qui exécute un programme laisse une trace
dans la mémoire, dans les journaux du système et souvent dans des fichiers
de préchargement~; qui copie des données laisse un volume sortant dans la
capture du lien Internet. Mais l'échange joue aussi dans l'autre sens~:
l'enquêteur qui examine une machine allumée y laisse ses propres traces. Le
principe d'échange oblige donc l'enquêteur autant que l'attaquant.

\subsection{Ce qui distingue la trace numérique}

\begin{tableaucours}{L{3.6cm}L{8.6cm}}{%
    \tableauHeader{Propriété} &
    \tableauHeader{Conséquence pour l'enquête}
}
Copiable à l'identique & On travaille sur une copie, et l'empreinte
cryptographique prouve que la copie égale l'original \\
Modifiable sans marque visible & Une date ou un journal se changent sans
rature~; seule la préservation prouve l'intégrité \\
Volatile & La mémoire vive et les connexions disparaissent à l'extinction~;
l'ordre de recueil compte \\
Abondante & Un serveur contient des millions de traces~; il faut des
questions précises pour savoir quoi chercher \\
Dépendante de son contexte & Une heure n'a de sens qu'avec son fuseau et la
dérive de l'horloge qui l'a produite \\
\end{tableaucours}
\vspace{-0.8em}
\begin{center}{\small\itshape Tableau~1.1. Propriétés de la trace numérique}\end{center}

% ============================================================================
\section{Les Principes d'un Examen Acceptable}
\label{sec:c01-principes}
% ============================================================================

La norme ISO/IEC~27037~\cite{iso27037_2012} fixe des lignes directrices pour
identifier, recueillir, acquérir et préserver la preuve numérique. Elle
retient quatre principes, que l'on paraphrase ici.

\begin{tableaucours}{L{3.2cm}L{5.4cm}L{3.6cm}}{%
    \tableauHeader{Principe} &
    \tableauHeader{Sens} &
    \tableauHeader{Comment le prouver}
}
Traçabilité (\textit{auditability}) & Un tiers peut reconstituer chaque geste
de l'enquêteur & Registre des actions, notes horodatées \\
Répétabilité & La même personne, avec les mêmes outils sur la même pièce,
obtient le même résultat & Empreintes identiques, versions d'outils notées \\
Reproductibilité & Une autre personne, avec d'autres outils, obtient un
résultat compatible & Méthode décrite, pièces conservées \\
Justification & Chaque choix de méthode ou d'outil peut être expliqué &
Raisons écrites dans le dossier \\
\end{tableaucours}
\vspace{-0.8em}
\begin{center}{\small\itshape Tableau~1.2. Quatre principes (d'après l'ISO/IEC~27037)}\end{center}

À ces quatre principes s'ajoutent deux règles de conduite que l'on retrouve
dans les guides de bonnes pratiques policiers, comme celui de l'ACPO
britannique~\cite{acpo2012}. La première est l'intervention minimale~: on ne
modifie pas une pièce, et si l'on y est contraint, par exemple pour capturer
la mémoire d'une machine allumée, on le fait en personne compétente et l'on
explique pourquoi. La seconde est la proportionnalité~: on recueille ce qui
répond aux questions de l'enquête, pas tout ce qui est techniquement
accessible.

\begin{boxinfo}
\textbf{Répétable, reproductible.} La distinction n'est pas un détail. Un
résultat répétable mais non reproductible dépend peut-être d'une
particularité de l'outil. Quand deux outils indépendants donnent la même
liste de fichiers supprimés, la conclusion gagne en confiance~; quand ils
divergent, la divergence elle-même devient un objet d'examen.
\end{boxinfo}

% ============================================================================
\section{Les Phases d'une Enquête}
\label{sec:c01-phases}
% ============================================================================

\subsection{Quatre phases selon le NIST}

Le guide NIST SP~800-86~\cite{nist80086} décrit l'intégration des techniques
forensiques à la réponse aux incidents en quatre phases.

\begin{center}
\begin{tikzpicture}[
    etape/.style = {draw=CouleurBordDef, fill=CouleurDef, rounded corners=3pt,
        minimum width=2.6cm, minimum height=1.0cm, align=center,
        font=\small\bfseries\color{CouleurPrimaire}},
    arrow/.style = {-{Stealth[length=4pt]}, thick, color=CouleurSecondaire}
]
\node[etape] (a) at (0,0)    {Collecte};
\node[etape] (b) at (3.3,0)  {Examen};
\node[etape] (c) at (6.6,0)  {Analyse};
\node[etape] (d) at (9.9,0)  {Rapport};
\draw[arrow] (a) -- (b); \draw[arrow] (b) -- (c); \draw[arrow] (c) -- (d);
\node[font=\footnotesize, align=center, text width=2.6cm] at (0,-1.1) {Support, données, intégrité};
\node[font=\footnotesize, align=center, text width=2.6cm] at (3.3,-1.1) {Extraire ce qui est pertinent};
\node[font=\footnotesize, align=center, text width=2.6cm] at (6.6,-1.1) {Relier, dater, conclure};
\node[font=\footnotesize, align=center, text width=2.6cm] at (9.9,-1.1) {Rendre compte, recommander};
\end{tikzpicture}
\end{center}
\begin{center}{\small\itshape Figure~1.1. Les quatre phases du NIST SP~800-86}\end{center}

La collecte identifie et recueille les supports et les données en préservant
leur intégrité. L'examen extrait, de la masse recueillie, les éléments
pertinents~: on décode un système de fichiers, on filtre des journaux. L'analyse
relie ces éléments entre eux pour répondre aux questions posées. Le rapport
en rend compte et formule des recommandations.

\subsection{Cinq classes de processus selon l'ISO/IEC 27043}

La norme ISO/IEC~27043~\cite{iso27043_2015} propose un modèle plus large,
qui commence avant l'incident. On en retient cinq classes de processus,
paraphrasées.

\begin{tableaucours}{L{3.4cm}L{8.8cm}}{%
    \tableauHeader{Classe} &
    \tableauHeader{Ce qu'elle couvre}
}
Préparation (\textit{readiness}) & Ce que l'organisation fait avant tout
incident pour qu'une enquête soit possible~: journaux conservés, horloges
synchronisées, procédure de première intervention \\
Initialisation & Le déclenchement~: détection, première évaluation,
décision d'enquêter, mandat \\
Acquisition & L'identification, le recueil, l'acquisition, le transport et
le stockage des pièces \\
Investigation & L'examen, l'analyse, l'interprétation, le rapport, la
présentation et la clôture \\
Processus concurrents & Ce qui accompagne toute l'enquête~: documentation,
chaîne de possession, préservation, autorisations \\
\end{tableaucours}
\vspace{-0.8em}
\begin{center}{\small\itshape Tableau~1.3. Classes de processus (d'après l'ISO/IEC~27043)}\end{center}

\begin{boxscenario}{Opération NYANGA. La préparation qui manquait}
\textbf{Constatations.} La sonde réseau installée en 2024 conserve 48~heures
de capture~; les journaux de la passerelle sont gardés 60~jours~; les
horloges des serveurs sont synchronisées sur un serveur de temps interne~;
l'établissement n'a aucune procédure écrite de première intervention, et
aucun journal des postes d'agence n'est centralisé.

\textbf{Analyse} (confiance élevée). La préparation existait pour le siège,
pas pour les agences. Si l'attaquant était passé par un poste d'agence (H2),
les traces auraient été dispersées sur vingt-deux sites et peut-être déjà
écrasées. La réaction de Mme~NDOUMBE a compensé l'absence de procédure~; un
autre responsable, ce vendredi-là, aurait pu laisser le DSI redémarrer.
\end{boxscenario}

% ============================================================================
\section{Trois Cadres d'Enquête}
\label{sec:c01-cadres}
% ============================================================================

Les mêmes pièces peuvent servir trois enquêtes différentes, qui n'ont ni les
mêmes questions, ni les mêmes exigences, ni les mêmes destinataires.

\begin{tableaucours}{L{2.6cm}L{3.2cm}L{3.2cm}L{3.2cm}}{%
    \tableauHeader{} &
    \tableauHeader{Enquête judiciaire} &
    \tableauHeader{Enquête interne} &
    \tableauHeader{Réponse à incident}
}
Qui la mandate & Le procureur, le juge, l'OPJ & La direction de
l'organisation & L'organisation victime \\
Question & Une infraction a-t-elle été commise, par qui~? & Un salarié
a-t-il manqué à ses obligations~? & Que s'est-il passé, comment arrêter et
reprendre~? \\
Exigence de preuve & La plus haute~: débat contradictoire & Celle du
contentieux du travail & Suffisante pour décider vite \\
Priorité & La vérité judiciaire & La décision disciplinaire & La reprise
de l'activité \\
Risque principal & Pièce écartée pour vice de forme & Atteinte à la vie
privée du salarié & Preuve détruite par la remise en route \\
\end{tableaucours}
\vspace{-0.8em}
\begin{center}{\small\itshape Tableau~1.4. Trois cadres d'enquête}\end{center}

La réponse à incident, décrite par le guide NIST SP~800-61~\cite{nist80061r2},
a pour objet premier de contenir l'attaque et de rétablir le service. Elle
entre en tension avec la préservation~: remettre en route un serveur,
c'est souvent effacer ce qu'il contenait. La règle pratique est de toujours
travailler comme si l'enquête pouvait finir devant un juge. Un incident qui
paraît interne au matin du vendredi peut devenir une affaire pénale la semaine suivante,
comme à NYANGA CRÉDIT~; les pièces mal recueillies le vendredi ne seront pas
recevables la semaine suivante.

\begin{boxinfo}
\textbf{À NYANGA, les trois cadres coexistent.} SANKOFA EXPERTISE conduit la
réponse à incident pour l'établissement~; l'OPJ et l'expert réquisitionné
conduisent l'enquête judiciaire~; si H3 avait survécu, la direction aurait dû
ouvrir une enquête interne, avec les garanties dues au salarié mis en cause.
Le Chapitre~\ref{chap:normes-droit} précise ce que la loi permet à chacun.
\end{boxinfo}

% ============================================================================
\section{Raisonner par Hypothèses Concurrentes}
\label{sec:c01-hypotheses}
% ============================================================================

Richards Heuer, analyste du renseignement américain, a décrit une méthode
simple pour résister à la tentation de confirmer sa première idée~: l'analyse
des hypothèses concurrentes~\cite{heuer1999}. On liste toutes les hypothèses
plausibles, on liste les pièces, puis on remplit une matrice en notant, pour
chaque couple, si la pièce est compatible ou incompatible avec l'hypothèse.
On retient l'hypothèse qui a le moins d'incompatibilités, pas celle qui a le
plus de compatibilités. Une pièce compatible avec toutes les hypothèses n'aide
pas à choisir, quelle que soit son importance apparente.

\begin{boxscenario}{Opération NYANGA. La matrice du 29 mars}
\begin{tableaucours}{L{5.2cm}cccc}{%
    \tableauHeader{Pièce} &
    \tableauHeader{H1} & \tableauHeader{H2} & \tableauHeader{H3} &
    \tableauHeader{H4}
}
Virements avec la clé de l'établissement & C & C & C & C \\
Session VPN \texttt{svc\_backup} depuis 203.0.113.18, le 2~mars & C & I & I &
I \\
Aucun échec d'authentification avant cette session & C & N & N & N \\
Aucun courriel piégé dans les boîtes des agences sur la période & C & I & C &
C \\
Lecture du fichier de clé sur SRV-CPT01, le 12~mars, par \texttt{adm.kamga}
& C & C & C & I \\
\texttt{adm.kamga} utilisé depuis le contrôleur de domaine, jamais depuis
un poste du siège & C & I & I & N \\
\end{tableaucours}

\smallskip
{\footnotesize C~: compatible~; I~: incompatible~; N~: neutre (ne distingue
pas).}

\medskip
\textbf{Constatations.} H1 n'a aucune incompatibilité~; H2 en a trois, H3
deux, H4 deux. La première ligne, la plus spectaculaire, est compatible avec
tout et ne sert à rien pour choisir.

\textbf{Analyse} (confiance élevée). H1 est retenue~; H2, H3 et H4 sont
incompatibles avec au moins deux pièces. Mais H3 n'est réfutée que telle
qu'elle était formulée~: un employé du siège n'a pas agi depuis un poste du
siège. Cela ne dit pas qu'aucun employé n'a participé~: un employé aurait pu
utiliser la passerelle depuis l'extérieur, comme l'attaquant. Cette question
est distincte de celle du chemin d'entrée, et aucune pièce ne la tranche.
\end{boxscenario}

\begin{boxlecture}{Les quatre hypothèses au 29 mars}
\begin{tableaucours}{L{6.0cm}cL{4.4cm}}{%
    \tableauHeader{Hypothèse} &
    \tableauHeader{Cote} &
    \tableauHeader{Pièce décisive}
}
H1, entrée par la passerelle & \OUI{} & Session du 2~mars, puis usage de
\texttt{adm.kamga} depuis le contrôleur de domaine \\
H2, poste d'agence infecté & \NON{} & Aucune session interne, aucun courriel
piégé \\
H3, employé agissant depuis le siège & \NON{} & Aucune session de
\texttt{adm.kamga} depuis un poste du siège \\
H4, fuite chez le prestataire & \NON{} & Clé lue sur SRV-CPT01 avant les
virements \\
Participation d'un employé, quel que soit le chemin & \NEK{} & Aucune pièce
ne l'établit ni ne l'exclut \\
\end{tableaucours}

\medskip
\textbf{Lecture.} Une proposition établie, trois réfutées, une non
tranchée. Le rapport intermédiaire ne doit pas dire « l'attaque vient de
l'extérieur » sans réserve~: il doit dire que l'entrée par la passerelle est
établie, et que la participation d'un employé n'est ni établie ni exclue.
\end{boxlecture}

% ============================================================================
\section{Les Biais de l'Enquêteur}
\label{sec:c01-biais}
% ============================================================================

L'enquêteur est un instrument de mesure, et cet instrument a des défauts
connus.

\begin{tableaucours}{L{3.4cm}L{4.6cm}L{4.2cm}}{%
    \tableauHeader{Biais} &
    \tableauHeader{Manifestation} &
    \tableauHeader{Parade}
}
Confirmation & On ne cherche que ce qui confirme l'hypothèse préférée &
Matrice des hypothèses concurrentes \\
Ancrage & La première information reçue pèse trop & Noter l'hypothèse
initiale et la réexaminer à chaque pièce \\
Autorité & L'avis du supérieur ou du client devient un fait & Exiger une
source pour chaque constatation \\
Contexte & Savoir qui est suspect oriente la lecture des traces & Examiner
les pièces avant de lire le dossier, quand c'est possible \\
Outil & Ce que l'outil n'affiche pas est tenu pour inexistant & Contrôler
avec un second outil, noter les limites \\
\end{tableaucours}
\vspace{-0.8em}
\begin{center}{\small\itshape Tableau~1.5. Biais et parades}\end{center}

Le biais d'outil mérite une attention particulière. Un logiciel d'analyse
affiche ce qu'il sait décoder. Un format de journal qu'il ne connaît pas, un
système de fichiers qu'il lit mal, une partition qu'il n'a pas montée
n'apparaissent nulle part, et l'absence d'affichage passe pour une absence de
trace. L'enquêteur qui écrit « aucune trace de connexion » doit pouvoir dire
« dans quelles sources, lues avec quels outils ». Sinon, la bonne cote est
\NEK{}, pas \NON{}.

% ============================================================================
\section{Ce que l'Investigation demande à chaque Métier}
\label{sec:c01-metiers}
% ============================================================================

\begin{boxmetier}{Analyste en réponse à incident}
\textbf{Ouvrir une enquête}
\begin{enumerate}
    \item Écrire les questions de l'enquête avant de toucher aux pièces.
    \item Lister toutes les hypothèses compatibles avec les premières
          constatations.
    \item Identifier, pour chacune, la pièce qui la réfuterait.
    \item Préserver ces pièces en priorité.
    \item Tenir le registre des actions dès la première minute.
\end{enumerate}
\end{boxmetier}

\begin{boxmetier}{Application OPJ/Expert Judiciaire}
L'expert ne choisit pas ses questions~: elles figurent dans la réquisition
ou l'ordonnance qui le désigne. S'il découvre un fait hors de sa mission, il
le signale à l'autorité qui l'a désigné, sans l'instruire de lui-même.
\end{boxmetier}

\begin{boxmetier}{RSSI}
La préparation est son affaire~: durée de conservation des journaux,
synchronisation des horloges, procédure de première intervention connue de
l'équipe de permanence. À NYANGA, il manquait la dernière.
\end{boxmetier}

\separateur

\begin{boxresume}
\textbf{Ce qu'il faut retenir de ce chapitre~:}
\begin{itemize}
    \item Une trace devient preuve quand on sait d'où elle vient, qu'on peut
          montrer qu'elle n'a pas changé, et qu'on peut expliquer ce qu'elle
          signifie.
    \item Quatre principes rendent un examen acceptable~: traçabilité,
          répétabilité, reproductibilité, justification.
    \item Le NIST décrit quatre phases (collecte, examen, analyse, rapport)~;
          l'ISO/IEC~27043 ajoute la préparation, avant l'incident.
    \item Enquête judiciaire, enquête interne et réponse à incident n'ont ni
          les mêmes questions ni les mêmes exigences~; on travaille toujours
          comme si l'affaire pouvait aller devant un juge.
    \item À NYANGA, la matrice du 29~mars établit H1, réfute H2, H3 et H4,
          et laisse \NEK{} la participation éventuelle d'un employé.
\end{itemize}
\end{boxresume}

\begin{boxinfo}
\textbf{Points de contrôle.} Vous avez compris ce chapitre si vous savez~:
\begin{itemize}
    \item expliquer la différence entre répétabilité et reproductibilité~;
    \item placer une action d'enquête dans l'une des phases du NIST~;
    \item remplir une matrice d'hypothèses concurrentes et en tirer une
          cote pour chaque hypothèse.
\end{itemize}
\end{boxinfo}

% ============================================================================
\section*{Exercices}
% ============================================================================

\begin{boxexo}[title={\ding{45}\quad Exercice~1.1. Trace ou preuve \niveaubase}]
Dites, pour chaque élément, s'il s'agit d'une simple trace ou d'une preuve
au sens du chapitre, et ce qui lui manque le cas échéant.
\begin{enumerate}[(a)]
    \item Une capture d'écran du journal de la passerelle, envoyée par
          WhatsApp par un technicien.
    \item L'export du même journal, sur support neuf, avec son empreinte
          SHA-256 inscrite au registre.
    \item Un fichier \texttt{lsass.dmp} trouvé sur SRV-AD01, dont personne
          n'a noté ni l'heure ni l'emplacement exacts.
    \item La capture mémoire de SRV-AD01 décrite au Chapitre~\ref{chap:prologue}.
\end{enumerate}
\end{boxexo}

\begin{boxexo}[title={\ding{45}\quad Exercice~1.2. Les phases et les classes \niveaubase}]
Rattachez chaque action à une phase du NIST SP~800-86 et à une classe de
processus de l'ISO/IEC~27043~: (a) synchroniser les horloges des serveurs
en 2024~; (b) débrancher le serveur de fichiers le 14~mars~; (c) extraire
du journal de la passerelle les sessions de \texttt{svc\_backup}~; (d)
relier ces sessions aux connexions de \texttt{adm.kamga}~; (e) remettre le
rapport à la direction.
\end{boxexo}

\begin{boxexo}[title={\ding{45}\quad Exercice~1.3. Une pièce qui change tout \niveaumoyen}]
On découvre le 2~avril que \texttt{adm.kamga} a ouvert, le 11~mars, une
session depuis un poste du siège, à 14~h~20, pendant les heures de bureau.
\begin{enumerate}[(a)]
    \item Ajoutez cette pièce à la matrice du 29~mars.
    \item Les cotes des hypothèses changent-elles~? Justifiez.
    \item Quelle vérification feriez-vous avant de modifier une cote~?
\end{enumerate}
\end{boxexo}

\begin{boxexo}[title={\ding{45}\quad Exercice~1.4. Trois cadres pour une affaire \niveauexpert}]
Six mois après l'incident, NYANGA CRÉDIT découvre qu'un agent comptable a
consulté, sans raison de service, des dossiers de clients figurant parmi les
fichiers exfiltrés.
\begin{enumerate}[(a)]
    \item Quels cadres d'enquête peuvent s'ouvrir~? Qui mandate chacun~?
    \item Pour chacun, formulez la question de fait et l'exigence de preuve.
    \item Quelles précautions prendre pour que les pièces recueillies dans
          l'enquête interne restent utilisables dans une enquête judiciaire~?
\end{enumerate}
\end{boxexo}

\bigskip

\begin{boxinfo}
\textbf{Pour aller plus loin.}
\begin{itemize}
    \item E.~Casey, \textit{Digital Evidence and Computer
          Crime}~\cite{casey2004}, pour les fondements.
    \item Le NIST SP~800-86~\cite{nist80086} et l'ISO/IEC~27043~\cite{iso27043_2015}.
    \item R.~Heuer, \textit{Psychology of Intelligence
          Analysis}~\cite{heuer1999}, chapitre sur les hypothèses
          concurrentes.
\end{itemize}
\end{boxinfo}

\begin{boxinfo}
\textbf{Transition.} Une hypothèse se tranche avec des pièces, et une pièce
ne vaut que si l'on peut prouver qu'elle n'a pas changé. C'est l'objet du
Chapitre~\ref{chap:preuve}~: la preuve numérique et la chaîne de possession.
\end{boxinfo}
